\section{The scaling toolbox}\label{sec:toolbox}

Throughout this section $U\in\R^{n\times d}$ is Parseval, $P=UU^T$,
$p=\diag P$, and $L_P=\Diag(p)-P\circ P$. For nonnegative symmetric weights
$w_{ij}$ on $[n]$ the weighted Laplacian $L$ acts by
$(Lx)_i=\sum_jw_{ij}(x_i-x_j)$; diagonal weights do not contribute. Unless
another Laplacian is specified, $L=L_P$. We write $\osc(x)=\max_ix_i-\min_ix_i$.

\subsection{Projections, distances, and the squared-Gram Laplacian}

\begin{lemma}[Frames versus projections]\label{lem:align}
Let $U,W\in\R^{n\times d}$ be Parseval with projections $P,Q$, and let
$\sigma_1,\dots,\sigma_d\in[0,1]$ be the singular values of $U^TW$. Then
\[
 \min_{R\in O(d)}\fro{U-WR}^2=2\sum_k(1-\sigma_k)\le
 \fro{P-Q}^2=2\sum_k(1-\sigma_k^2)=2\fro{(I-P)W}^2\le2\fro{U-W}^2 .
\]
Right multiplication by $R\in O(d)$ preserves Parsevalness and every row norm.
If $X\in\R^{n\times d}$ has $\delta=\opn{X^TX-I}<1$, its polar factor
$\tilde X=X(X^TX)^{-1/2}$ is Parseval, $\fro{X-\tilde X}^2\le d\delta^2$, and
$\norm{x_i}^2/(1+\delta)\le\norm{\tilde x_i}^2\le\norm{x_i}^2/(1-\delta)$. If
moreover $\delta\le\frac12$, then $|\norm{\tilde x_i}^2-\norm{x_i}^2|\le2\delta\norm{x_i}^2$
and row $i$ of $\tilde X\tilde X^T-XX^T$ has squared norm at most
$6\delta^2\norm{x_i}^2$.
\end{lemma}

\begin{proof}
$\fro{U-WR}^2=2d-2\tr(U^TWR)$, maximised over $R$ by a polar factor of $U^TW$ (any orthogonal $R$ with $U^TWR\succeq0$; it need not be unique),
giving $2d-2\sum_k\sigma_k$. Next $\fro{P-Q}^2=2d-2\tr(PQ)=2d-2\fro{U^TW}^2$
and $\fro{(I-P)W}^2=d-\fro{U^TW}^2$; also $1-\sigma\le1-\sigma^2$ on $[0,1]$,
and $(I-P)W=(I-P)(W-U)$ gives $\fro{(I-P)W}\le\fro{U-W}$. For the polar factor, if $\varsigma_k$ are the singular values of $X$
then $\fro{X-\tilde X}^2=\sum_k(\varsigma_k-1)^2\le\sum_k(\varsigma_k^2-1)^2$, and
$\norm{\tilde x_i}^2=x_i^T(X^TX)^{-1}x_i$. With $E=(X^TX)^{-1}-I$, so
$\opn E\le\delta/(1-\delta)\le2\delta$, row $i$ of $\tilde X\tilde X^T-XX^T=XEX^T$
is $XEx_i$, of squared norm at most $(1+\delta)4\delta^2\norm{x_i}^2$.
\end{proof}

\begin{lemma}[Energy identity]\label{lem:energy}
$L$ is the Laplacian of the weights $P_{ij}^2$ $(i\ne j)$, $L\succeq0$,
$L\one=0$, $L_{I-P}=L_P$, and for $x\in\R^n$, $X=\Diag(x)$,
\[
 x^TLx=\tfrac12\sum_{i,j}P_{ij}^2(x_i-x_j)^2=\fro{(I-P)XU}^2
 =\tfrac12\fro{[X,P]}^2 .
\]
\end{lemma}

\begin{proof}
$\sum_jP_{ij}^2=(P^2)_{ii}=p_i$, so the rows of $\Diag(p)-P\circ P$ sum to zero
and its off-diagonal entries are $-P_{ij}^2$. Next
$\fro{(I-P)XU}^2=\tr(XPX(I-P))=\sum_ix_i^2p_i-\sum_{i,j}x_ix_jP_{ij}^2$, and
$[X,P]_{ij}=(x_i-x_j)P_{ij}$. For $I-P$ the off-diagonal weights are again
$P_{ij}^2$.
\end{proof}

\begin{lemma}[Scaled projections]\label{lem:scaled}
For $w\in\R^n_{>0}$ put $D=\Diag(w)$, $M=(U^TD^2U)^{-1/2}$ and $W=DUM$. Then
$W$ is Parseval with column space $D\im U$, and
\[
 \fro{WW^T-P}^2=2\fro{(I-P)W}^2\le\frac{2\,w^TLw}{\min_iw_i^2} .
\]
\end{lemma}

\begin{proof}
$W^TW=MU^TD^2UM=I$ and $U^TD^2U\succeq(\min_iw_i)^2I$, so
$\opn M\le1/\min_iw_i$. By \cref{lem:align,lem:energy},
$\fro{(I-P)W}^2\le\opn M^2\fro{(I-P)DU}^2=\opn M^2\,w^TLw$.
\end{proof}

\subsection{Diagonal scaling}

For $s\in\R^n$ let $D_s=\Diag(e^{s})$, let $P(s)$ be the projection onto
$D_s\im U$, $r(s)=\diag P(s)$, and $F(s)=\frac12\log\det(U^TD_s^2U)$.

\begin{lemma}\label{lem:potential}
$F$ is smooth and convex, $\nabla F(s)=r(s)$, $\sum_ir_i(s)=d$, and
$F(s+c\one)=F(s)+cd$.
\end{lemma}

\begin{proof}
With $G=U^TD_s^2U=\sum_ie^{2s_i}u_iu_i^T$ one has
$\partial_iF=e^{2s_i}u_i^TG^{-1}u_i=P(s)_{ii}$, because $P(s)=D_sUG^{-1}U^TD_s$;
differentiating again gives $\nabla^2F=2L_{P(s)}\succeq0$. The trace of $P(s)$ is
$d$, and $G$ scales by $e^{2c}$ under $s\mapsto s+c\one$.
\end{proof}

\begin{lemma}[One-sided scaling inequalities]\label{lem:scaling-ineq}
For $z\in\R^n_{>0}$ let $r$ be the diagonal of the projection onto
$\Diag(\sqrt z)\im U$. Then, coordinatewise,
\[
 Lz\le z\circ(r-p),\qquad L(z^{-1})\le-z^{-1}\circ(r-p).
\]
\end{lemma}

\begin{proof}
Let $G=U^T\Diag(z)U$. Since $G+z_i^2G^{-1}-2z_iI=(G^{1/2}-z_iG^{-1/2})^2\succeq0$,
evaluating at $u_i$ and using $u_i^TGu_i=\sum_jP_{ij}^2z_j=p_iz_i-(Lz)_i$ and
$z_i^2u_i^TG^{-1}u_i=z_ir_i$ gives the first inequality. Let
$V\in\R^{n\times(n-d)}$ be an isometry onto $(\im U)^\perp$. The subspaces
$\Diag(\sqrt z)\im U$ and $\Diag(z^{-1/2})\im V$ are orthogonal complements, so
the projection onto the latter has diagonal $\one-r$, while $VV^T=I-P$ has
diagonal $\one-p$ and Laplacian $L$ (\cref{lem:energy}). The first inequality
for $(V,z^{-1})$ is the second. (If $d=n$ then $L=0$ and $r=p=\one$.)
\end{proof}

\subsection{The barrier constant and static balancing}

\begin{definition}[Barrier constant]\label{def:barrier}
Let $L$ be a weighted Laplacian on $[n]$. For $S\subseteq[n]$, an
\emph{$S$-barrier} is a vector $\psi\ge0$ with $(L\psi)_i\ge1$ for every
$i\notin S$. The \emph{barrier constant} $H(L)\in(0,\infty]$ is the least $H$ such
that every $S$ with $|S|\ge n/2$ has an $S$-barrier with $\norm\psi_\infty\le H$
($H(L)=\infty$ if some such $S$ has no barrier).
\end{definition}

For the continuous-time random walk that jumps from $i$ to $j$ at rate
$w_{ij}$, the expected hitting time $h_S(x)=\E_x\tau_S$ of $S$ vanishes on $S$
and satisfies $(Lh_S)_i=1$ off $S$. If $\psi$ is any $S$-barrier and $\beta'>1$,
the maximum principle argument in the proof of \cref{lem:median}, applied to
$h_S-\beta'\psi$, gives $h_S\le\beta'\psi$. So $h_S$ is the minimal $S$-barrier,
and $H(L_P)$ is the half-set hitting time of the walk with rates $P_{ij}^2$: the
largest expected time, over all starting points and all $S$ with $|S|\ge n/2$,
needed to reach $S$. If $H(L)<\infty$ then the graph is connected: for a
component $J$ with $|J|\le n/2$, summing $(L\psi)_i$ over $i\in J$ gives $0$, so
$S=[n]\setminus J$ has no barrier. The barrier constant is within a factor $2$ of
the $\ell_\infty$ Poisson constant of $L$ (\cref{rem:poisson}). A small barrier
constant controls scalings in $\ell_\infty$:

\begin{lemma}[Median barrier]\label{lem:median}
Let $L$ be a weighted Laplacian on $[n]$ with $H(L)\le H<\infty$, let $\beta\ge0$ with $\beta H<1$, let
$z\in\R^n_{>0}$, and let $m$ be a median of $z$ (so that $\{z\le m\}$ and
$\{z\ge m\}$ both have at least $n/2$ elements). If
\[
 (Lz)_i\le\beta z_i\ \text{ when }z_i>m,\qquad
 (Lz^{-1})_i\le\beta z_i^{-1}\ \text{ when }z_i<m ,
\]
then $(1-\beta H)\max_iz_i\le m\le(1-\beta H)^{-1}\min_iz_i$.
\end{lemma}

\begin{proof}
Let $S=\{z\le m\}$, let $\psi$ be an $S$-barrier with $\norm\psi_\infty\le H$, let
$M=\max z$ and $\beta'>\beta$. The vector $v=z-m\one-\beta'M\psi$ is $\le0$ on
$S$, and for $i\notin S$, $(Lv)_i\le\beta z_i-\beta'M<0$. At a maximiser $i$ of
$v$, $(Lv)_i=\sum_jw_{ij}(v_i-v_j)\ge0$; so every maximiser lies in $S$, and
$v\le0$. Hence $M\le m+\beta'MH$; let $\beta'\downarrow\beta$. Apply the same
argument to $z^{-1}$, whose median is $m^{-1}$, with $S=\{z\ge m\}$.
\end{proof}

\begin{theorem}[Static balancing]\label{thm:balancing}
Suppose $H(L_P)\le H<\infty$, and $q\in\R^n$ satisfies $\sum_iq_i=d$ and
$\beta H\le\frac12$, where $\beta=\norm{q-p}_\infty$. Then
there is a Parseval $W$ with $\diag(WW^T)=q$ and
\[
 \fro{U-W}^2\le\fro{UU^T-WW^T}^2\le\tfrac{15}4\norm{q-p}_1 .
\]
$W$ is obtained from $U$ by a row scaling with ratio of scales at most $2$,
whitening, and an orthogonal right factor.
\end{theorem}

\begin{proof}
Let $s$ minimise $\Phi(s)=F(s)-\ip qs$ over the cube $[0,1]^n$ (compactness),
and write $r=r(s)$. Since $\partial_i\Phi=r_i-q_i$ (\cref{lem:potential}),
coordinatewise optimality gives
\begin{equation}\label{eq:kkt}
 r_i\le q_i\ \text{ if }s_i>0,\qquad r_i\ge q_i\ \text{ if }s_i<1 .
\end{equation}
Let $z=e^{2s}$ with median $m$, so $1\le\min z\le m\le\max z\le e^{2}$. If
$z_i>m$ then $s_i>0$, so by \cref{lem:scaling-ineq} and~\eqref{eq:kkt}
$(Lz)_i\le z_i(r_i-p_i)\le z_i(q_i-p_i)\le\beta z_i$; if $z_i<m$ then $s_i<1$
and likewise $(Lz^{-1})_i\le\beta z_i^{-1}$. By \cref{lem:median},
$\max z/\min z\le(1-\beta H)^{-2}\le4$, i.e.\ $\osc(s)\le\log2<1$. Hence
$s$ does not meet both faces of the cube: either all $s_i>0$, and then
$r\le q$, or some $s_j=0$, all $s_i<1$, and then $r\ge q$. Since
$\sum r_i=d=\sum q_i$, in both cases $r=q$.

For the distance let $w=e^{s}$, $f=q-p$, and let $c$ be the midpoint of the
range of $z^2$. \Cref{lem:scaling-ineq} with $r=q$ gives $Lz\le z\circ f$, so,
using $\sum_if_i=0$,
\[
 z^TLz\le\sum_iz_i^2f_i=\sum_i(z_i^2-c)f_i\le\tfrac12\osc(z^2)\norm f_1 .
\]
As $(z_i-z_j)^2\ge4\min w^2\,(w_i-w_j)^2$, \cref{lem:energy,lem:scaled} give
\[
 \fro{P(s)-P}^2\le\frac{2\,w^TLw}{\min w^2}\le\frac{z^TLz}{2\min z^2}
 \le\frac{\osc(z^2)}{4\min z^2}\norm f_1\le\frac{4^2-1}4\norm f_1 .
\]
Finally apply \cref{lem:align} to the whitened frame $DU(U^TD^2U)^{-1/2}$,
$D=\Diag(w)$, whose projection is $P(s)$.
\end{proof}

\begin{remark}
The hypothesis forces $q_i\ge p_i/2>0$: for $n\ge2$, the barrier for
$S=[n]\setminus\{i\}$ gives $1\le(L\psi)_i\le p_i(1-p_i)H$, so
$\beta\le1/(2H)\le p_i/2$. The proof uses no existence theorem for scalings,
no flow, and no genericity or connectivity argument: the box supplies a
minimiser, the median barrier keeps it away from one face, and the trace
identity $\sum_ir_i=d$ does the rest. With
$\beta H<1$ and a box of side $T>-\log(1-\beta H)$ the same proof gives scale
ratio $(1-\beta H)^{-1}$ and cost $\frac14((1-\beta H)^{-4}-1)\norm{q-p}_1$.
\end{remark}

\subsection{Barriers from a dense core}

\begin{lemma}[Dense core with exceptional vertices]\label{lem:core}
Let $L$ be the Laplacian of nonnegative symmetric weights $w_{ij}$ on $[n]$,
and let $[n]=\mathcal G\sqcup\mathcal B$. Suppose:
\begin{enumerate}[label=(\alph*)]
\item there is $\kappa>0$ with $\sum_{j\in S}w_{ij}\ge\kappa$ for every
$i\in\mathcal G$ and every $S\subseteq[n]\setminus\{i\}$ with $|S|\ge n/2$;
\item there are $y\in\R^{\mathcal B}$, $y\ge0$, and $R,F_{\mathcal B}\ge0$ with
$\norm y_\infty\le R$, $(L_{\mathcal B\mathcal B}y)_i\ge1$ for $i\in\mathcal B$,
and $\sum_{j\in\mathcal B}w_{ij}y_j\le F_{\mathcal B}$ for $i\in\mathcal G$.
(If $\mathcal B=\varnothing$ take $R=F_{\mathcal B}=0$.)
\end{enumerate}
Then $H(L)\le(1+F_{\mathcal B})/\kappa+R$. Moreover, with $b=|\mathcal B|$:
\begin{enumerate}[label=(\roman*)]
\item (a) holds with $\kappa=\vartheta\gamma$ if $\vartheta\in(0,\frac12]$,
$\gamma>0$, and every $i\in\mathcal G$ has at least $(\frac12+\vartheta)n$ indices
$j\ne i$ with $w_{ij}\ge\gamma/n$;
\item (b) holds with $R=\sqrt b/\lambda$, $F_{\mathcal B}=b$, if
$x^TLx\ge\lambda\norm x^2$ for every $x$ supported on $\mathcal B$ ($\lambda>0$);
\item for $w_{ij}=P_{ij}^2$, (b) holds with $R=\frac8{3\alpha}$,
$F_{\mathcal B}=\frac2{3\alpha}\max_ip_i$ if $\alpha>0$, $p_i\ge\alpha/2$ on
$\mathcal B$ and $\sum_{i\in\mathcal B}p_i\le\frac14$.
\end{enumerate}
\end{lemma}

\begin{proof}
Let $|S|\ge n/2$, put $\lambda_0=(1+F_{\mathcal B})/\kappa$, extend $y$ by $0$ to
$\mathcal G$, and let $\psi=(\lambda_0\one+y)\circ\one_{S^c}$, so
$0\le\psi\le\lambda_0+R$ and $\psi_j\le\lambda_0+y_j$ for all $j$. For
$i\in\mathcal G\setminus S$, $\psi_i=\lambda_0$ and $\psi_j=0$ on $S$, so
\[
 (L\psi)_i\ge\lambda_0\sum_{j\in S}w_{ij}-\sum_{j\in\mathcal B}w_{ij}y_j\ge
 \lambda_0\kappa-F_{\mathcal B}=1 .
\]
For $i\in\mathcal B\setminus S$, $\psi_i=\lambda_0+y_i$, so
$(L\psi)_i\ge\sum_jw_{ij}(y_i-y_j)=(L_{\mathcal B\mathcal B}y)_i\ge1$. Thus $\psi$
is an $S$-barrier.

(i) If $N_i$ is the set of such $j$ and $S\subseteq[n]\setminus\{i\}$, then
$|N_i\cap S|\ge|N_i|+|S|-(n-1)>\vartheta n$.
(ii) $\mathsf L=L_{\mathcal B\mathcal B}\succeq\lambda I$; let $y=\mathsf L^{-1}\one_{\mathcal B}$. If
$y$ had a negative minimum at $i$, then $(\mathsf Ly)_i=y_i\sum_{j\notin\mathcal
B}w_{ij}+\sum_{j\in\mathcal B}w_{ij}(y_i-y_j)\le0$, a contradiction; so $y\ge0$,
and $\norm y_\infty\le\norm y_2\le\sqrt b/\lambda$. Since
$(\mathsf L\one_{\mathcal B})_j=\sum_{i\in\mathcal G}w_{ij}$ for $j\in\mathcal B$, for each
$i_0\in\mathcal G$
\[
 \sum_{j\in\mathcal B}w_{i_0j}y_j\le\sum_{i\in\mathcal G}\sum_{j\in\mathcal B}w_{ij}y_j
 =\one_{\mathcal B}^T\mathsf Ly=b .
\]
(iii) $P_{ij}^2\le p_ip_j$ gives, for $i\in\mathcal B$,
$(L_{\mathcal B\mathcal B}\one_{\mathcal B})_i=p_i-\sum_{j\in\mathcal B}P_{ij}^2\ge
p_i(1-\sum_{j\in\mathcal B}p_j)\ge\frac{3\alpha}8$, and for $i\in\mathcal G$,
$\sum_{j\in\mathcal B}P_{ij}^2\le p_i/4$. Take $y=\frac8{3\alpha}\one_{\mathcal B}$.
\end{proof}

\begin{remark}[Poisson constants]\label{rem:poisson}
The $\ell_\infty$ Poisson constant $K$ of $L$ is the least $K$ such that every
$f\perp\one$ has a solution of $Lg=f$ with $\norm g_\infty\le K\norm f_\infty$. Then $K\le H\le2K$. Indeed, for
$|S|\ge n/2$ the forcing $f=\one-\frac n{|S|}\one_S$ has $\norm f_\infty\le1$,
and $g+K\one$ is an $S$-barrier; conversely, if $H<\infty$ the graph is connected,
so $Lg=f$ is solvable for every $f\perp\one$, and if $\norm f_\infty\le1$ and $m$ is
a median of $g$, the argument of \cref{lem:median}
applied to $g-m\one-\beta'\psi$ ($\beta'>1$) gives $g\le m+H$, and symmetrically
$g\ge m-H$. Barriers are more convenient than Poisson solutions because they
can be written down explicitly, as in \cref{lem:core}.
\end{remark}

\subsection{From a seed to an exact frame}

Both constructions produce the following object.

\begin{definition}[Seed]\label{def:seed}
Let $\Theta\ge2$ and $t>0$. A \emph{$(\Theta,t)$-seed} for $X\in\R^{n\times d}$
is a Parseval $V\in\R^{n\times d}$ with
\[
 \fro{V-X}^2\le\Theta t^2d,\qquad
 H(L_{VV^T})\le\frac{\Theta}{at^2},\qquad
 \max_i\bigl|(VV^T)_{ii}-a\bigr|\le\frac{at^2}{2\Theta}.
\]
\end{definition}

\begin{corollary}\label{cor:seed}
If $X$ has a $(\Theta,t)$-seed, there is an ENP frame $W$ with
$\fro{X-W}^2\le3\Theta t^2d$.
\end{corollary}

\begin{proof}
Apply \cref{thm:balancing} to $V$ with $q=a\one$: $\beta H\le\frac12$, and
$\fro{V-W}^2\le\frac{15}4n\cdot\frac{at^2}{2\Theta}\le t^2d$. Then
$\fro{X-W}^2\le2\fro{X-V}^2+2\fro{V-W}^2\le(2\Theta+2)t^2d\le3\Theta t^2d$.
\end{proof}
